% chapter: はじめに

\chapter{はじめに}\label{章：はじめに}
\begin{bookcolumns}
  \section{本書について}\label{節：はじめに}
    本書は、\bookterm{recurrence}{漸化式}の解法と、その背後にある考え方をたどる本である。
    差を取る。倍率をそろえる。複数の項を組み合わせる。
    一つの式変形から、別の解法や、グラフ・\bookterm{linear}{線形}代数へのつながりが見えてくる。
    その道筋を、式と短い文章で描いていく。

    各型では、代表的な解法から始める。
    例題の方針には着眼点を、解答には結果を、解法には計算の過程を記した。
    例題を解いた後で、同じ操作を別の見方から捉え直す。
    計算の理由と、複数の方法の使い分けも扱う。

    紙とペンで式を確かめると、文章だけでは見えにくい変化も見えてくる。
    演習は、その理解を別の式へ広げるためにある。
    本書でいう「基本」は、後の議論を支える考え方を指す。
    計算の道具は\chapterref{章：前提知識と計算の基本}へ、
    用語と発展的な計算は\chapterref{章：用語と計算の参照}へ戻って確認できる。

    \bookterm{general}{一般項}を求めることから、\bookterm{limit}{極限}や整数の性質を調べることへ。
    それぞれの解法を結び付けながら、\bookterm{recurrence}{漸化式}を見渡していく。

  \section{章立て}\label{節：章立て}
    基本の計算から、解法、その選び方、数学的なつながりへ進む。
    巻末には、言葉と計算の参照先を置いた。
    \begin{itemize}
      \item \chapterref{章：前提知識と計算の基本}は参照用である.必要な知識が出てきたときに,対応する節へ戻ればよい.
      \item \chapterref{章：漸化式の解説}では,個々の\bookterm{recurrence}{漸化式}に対してのアプローチを学ぶ.
      \item \chapterref{章：実際の解き方}では,\chapterref{章：漸化式の解説}で学んだ内容をまとめ,実際に\bookterm{recurrence}{漸化式}が与えられた時にどうするかを学ぶ.
      \item \chapterref{章：付録}では、本編の操作を\bookterm{matrix}{行列}・\bookterm{difference-operator}{差分}・\bookterm{generating}{母関数}から見直す。
      \item \chapterref{章：用語と計算の参照}には、用語の意味と、\bookterm{logarithm}{対数}・\bookterm{trigonometric}{三角関数}・\bookterm{limit}{極限}・複素数の計算をまとめた。
    \end{itemize}
    それぞれの関係は
    \BookFigRoadmap
    のようになっている.

  \section{対象読者と読む経路}\label{節：対象読者}
    対象読者は,\bookterm{recurrence}{漸化式}を初めて学ぶ人から,既習の解法の理由や数学的なつながりを知りたい人まで対応している.
    特にこの本にマッチするのは,数学が好きな高校生である.
    関心や既習の内容に合わせて、次の場所から読める。

    \noindent\textbf{【初めて学ぶ人】}\par
    \sectionref*{節：数列と漸化式}で記号の読み方を確かめ、
    \sectionref*{節：漸化式を解く前に}から\sectionref*{節：基本4パターン}へ進む。
    その後、\sectionref*{節：追う量を変える}、
    \sectionref*{節：解を組み合わせる見方},\sectionref*{節：応用7パターン}へ進む.
    一つの型を読んだら,その例題と\hyperref[節：演習の読む経路]{「演習の読む経路」}の代表題に取り組む.
    到達目標は,変形後の式を計算できることに加えて,「何に注目し,何を簡単にしたか」を一文で説明し,\bookterm{initial}{初期条件}と元の式への代入で答えを確かめられることである.

    \noindent\textbf{【一度学んだ人】}\par
    \sectionref*{節：解法の決定}と\sectionref*{節：操作を選ぶ練習}から入り,理由を説明できない操作だけ第3章へ戻る.
    \hyperref[節：演習の読む経路]{「演習の読む経路」}の比較題では,別解の手掛かり・計算量・条件の違いを調べる.
    到達目標は,同じ問題の二つの解法を比較し,係数や初項が変わったときにも方法を選び直せることである.

    \noindent\textbf{【数学を深く学びたい人】}\par
    発展パターン・特殊性質型と発展演習を,付録の関連する節と往復して読む.
    到達目標は,異なる型で使った操作の共通点を述べ,\bookterm{general}{一般項}・\bookterm{limit}{極限}・\bookterm{invariant}{不変量}など,問題に合う目標を選べることである.
    \bookterm{matrix}{行列}は付録の冒頭で導入する.微分を使う箇所では接線と導関数の知識が必要であり,\bookterm{trigonometric}{三角関数}・\bookterm{complex-plane}{複素平面}・\bookterm{limit}{極限}は、巻末の参照資料で確認できる。

  \section{本書でのルール}\label{節：本書でのルール}
    左罫線と薄い背景のある囲みは,読み返すための着眼点・条件・結論を示す.
    太字の小見出しで,発想・計算・解釈の切り替わりを示す.
    議論の道筋は、本文の文章と式に記した。
    用語には、巻末の説明へのリンクがある。
    下段の注釈には、条件の確認や短い補足を置いた。
    操作に必要な条件は、本文にも示す。
    PDFの目次・見出し参照・例題の見出しはリンクになっている.
    例題・演習問題の番号から解答へ,方針・解答・解法の番号から問題へ移動できる.

    特に断りが無い限りは,\,$k,l,m,n$は1以上の整数,\,$a,b,c,d,p,q,r,s,t,u,v,x,y,\lambda$は実数として扱う.
    また,\,$\alpha,\beta,\gamma,\delta$は方程式の解として扱う.
    ただし$d$は\bookterm{arithmetic}{等差数列}の公差を表す文字として,\,$0$を含む任意の実数を取りうる.また,添字は$n\ge1$を基本とするが,\,\sectionref{節：虚数解型隣接3項間漸化式}や付録の一部では$a_0$から始まる添字を用いることがある.
    また,\,$z$は複素数として扱う.
    また,\,$\therefore\,,\because\,,\Rightarrow\,,\Leftrightarrow$は左から順に「従って」「なぜならば」「左側が成り立つならば右側が成り立つ」「同値である」を意味する.
    加えて,\,$\cdots$や$\vdots$は途中を省略する時に用いる.

\end{bookcolumns}
